\chapter{Patterns: render budgets}
\label{ch:renderbudget}

\begin{aside}
For an app to feel snappy, each render must finish in time.
A budget per phase makes the target explicit.
\end{aside}

\section{Per-phase budget}

For a 60 Hz target (16.7 ms total):

\begin{itemize}[leftmargin=*]
  \item Layout: 1 ms.
  \item Paint: 2 ms.
  \item Diff: 1 ms.
  \item Emit + write: 1 ms.
  \item Application code: 11 ms.
\end{itemize}

Most apps have plenty of headroom.

\section{Measuring}

The trace facility (Chapter~\ref{ch:tracing}) gives per-phase
numbers. Compare against budget.

\section{Going over}

A frame that exceeds budget delays the next. Repeated
overruns produce visible lag.

\section{Recovery}

If a frame takes too long, the next frame fires
immediately (no wait). The loop catches up. Brief
overruns are invisible; sustained overruns are felt.

\section{Adaptive quality}

For complex apps, drop quality when over budget:

\begin{itemize}[leftmargin=*]
  \item Lower animation frame rate.
  \item Defer non-essential rendering.
  \item Simplify expensive widgets.
\end{itemize}

\section{Recap}

\begin{recap}
\begin{itemize}[leftmargin=*]
  \item Budget per phase.
  \item Trace and compare.
  \item Adaptive quality for overflow.
\end{itemize}
\end{recap}

\section*{Workshop}

\begin{enumerate}[leftmargin=*]
  \item Set budgets for your app's phases.
  \item Trace and identify overruns.
  \item Add adaptive quality where useful.
\end{enumerate}
